Kurs:Algebraische Kurven (Osnabrück 2025-2026)/Arbeitsblatt 9/latex
\setcounter{section}{9}






  \zwischenueberschrift{Übungsaufgaben}


\inputaufgabe
{}
{




Begründe, warum der Ring

\mathdisp {\Z[X,Y,Z,W]/ \left(  XY-ZW, 5X^8-YZ^3+2WXY  \right)} {  }

\definitionsverweis {noethersch}{}{}
ist.




}
{} {}


\inputaufgabe
{}
{




Es sei $R$ ein
\definitionsverweis {kommutativer Ring}{}{}
und sei

\mavergleichskettedisp
{\vergleichskette
{ {\mathfrak a}_1
}
{ \subseteq} { {\mathfrak a}_2
}
{ \subseteq} { {\mathfrak a}_3
}
{ \subseteq} { \ldots
}
{ } {
}
}
{}{}{}
eine aufsteigende Kette von
\definitionsverweis {Idealen}{}{.}
Zeige, dass die
\definitionsverweis {Vereinigung}{}{}

\mathl{\bigcup_{n \in \N} {\mathfrak a}_n}{} ebenfalls ein Ideal ist. Zeige durch ein einfaches Beispiel, dass die Vereinigung von Idealen im Allgemeinen kein Ideal sein muss.




}
{} {}


\inputaufgabe
{}
{




Zeige, dass das
\definitionsverweis {Produkt}{}{}
$R \times S$ zu
\definitionsverweis {noetherschen Ringen}{}{}
\mathkor {} {R} {und} {S} {}
wieder noethersch ist.




}
{} {}


\inputaufgabe
{}
{




Es sei $K$ ein
\definitionsverweis {Körper}{}{.} Zeige, dass es in
\mathl{K[X,Y]}{} keine obere Schranke für die Anzahl der Erzeuger von
\definitionsverweis {Idealen}{}{}
\zusatzklammer {in einem minimalen
\definitionsverweis {Erzeugendensystem}{}{}} {} {}
gibt.




}
{} {Tipp: Betrachte die Potenzen
\mathl{(X,Y)^m}{.}}


\inputaufgabe
{}
{




Es sei $K$ ein Körper und sei

\mathdisp {K[X_n, \, n \in \N]} {  }

der Polynomring über $K$ in unendlich vielen Variablen. Man beschreibe darin ein nicht endlich erzeugtes Ideal und eine unendliche, echt aufsteigende Idealkette.




}
{} {}


\inputaufgabegibtloesung
{}
{




Zeige, dass ein
\definitionsverweis {Unterring}{}{}

\mavergleichskette
{\vergleichskette
{R
}
{ \subseteq }{ S
}
{  }{
}
{    }{
}
{   }{
}
}
{}{}{}
eines
\definitionsverweis {noetherschen Ringes}{}{}
nicht noethersch sein muss.




}
{} {}


\inputaufgabe
{}
{




Man gebe ein Beispiel eines nicht-noetherschen Ringes, dessen \definitionsverweis {Reduktion}{}{}
ein Körper ist.




}
{} {}


\inputaufgabe
{}
{




Es sei $R$ ein
\definitionsverweis {kommutativer Ring}{}{}
und sei

\mavergleichskette
{\vergleichskette
{ {\mathfrak a}
}
{ \subset }{ R
}
{  }{
}
{    }{
}
{   }{
}
}
{}{}{}
ein
\zusatzklammer {echtes} {} {}
\definitionsverweis {Ideal}{}{}
mit dem
\definitionsverweis {Restklassenring}{}{}

\mathl{R/ {\mathfrak a}}{.} Zeige durch ein Beispiel, dass ${\mathfrak a}$ endlich erzeugt und

\mathl{R/ {\mathfrak a}}{}
\definitionsverweis {noethersch}{}{}
sein kann, ohne dass $R$ noethersch ist.




}
{} {}


\inputaufgabe
{}
{




Es sei $R$ ein
\definitionsverweis {kommutativer Ring}{}{}
und

\mavergleichskette
{\vergleichskette
{ {\mathfrak b}
}
{ \subseteq }{ R[X]
}
{  }{
}
{    }{
}
{   }{
}
}
{}{}{}
ein
\definitionsverweis {Ideal}{}{,}
das zumindest ein
\definitionsverweis {normiertes Polynom}{}{}
enthalte. Was bedeutet dies für die im Beweis zum

Hilbertschen Basissatz

konstruierte Idealkette in $R$?




}
{} {}


\inputaufgabe
{}
{




Es sei $R$ ein
\definitionsverweis {kommutativer Ring}{}{.}
Charakterisiere diejenigen
\definitionsverweis {Ideale}{}{}

\mavergleichskette
{\vergleichskette
{ {\mathfrak b}
}
{ \subseteq }{ R[X]
}
{  }{
}
{    }{
}
{   }{
}
}
{}{}{}
mit der Eigenschaft, dass die im Beweis zum

Hilbertschen Basissatz

konstruierte Idealkette in $R$ konstant ist.




}
{} {}


\inputaufgabe
{}
{




Es sei $K$ ein
\definitionsverweis {Körper}{}{} und sei

\mavergleichskette
{\vergleichskette
{ R
}
{ = }{ K[X]
}
{  }{
}
{    }{
}
{   }{
}
}
{}{}{}
der
\definitionsverweis {Polynomring}{}{}
über $K$. Bestimme zu den
\definitionsverweis {Idealen}{}{}

\mavergleichskettedisp
{\vergleichskette
{ {\mathfrak b}_m
}
{ =} { (X,Y)^m
}
{ \subseteq} { R[Y] = K[X,Y]
}
{ } {
}
{ } {
}
}
{}{}{}
die im Beweis zum

Hilbertschen Basissatz

konstruierte Idealkette in $R$. Wann wird sie stationär?




}
{} {}


\inputaufgabe
{}
{




Bestimme zum
\definitionsverweis {Ideal}{}{}

\mavergleichskettedisp
{\vergleichskette
{ I
}
{ =} {  \left(  6,6x^2+2x+3,3x^3+5,2x^5+x-4,4x^7-3x  \right)
}
{ } {
}
{ } {
}
{ } {
}
}
{}{}{}
in $\Z[x]$ die im Beweis zum

Hilbertschen Basissatz

konstruierte Idealkette und das zugehörige Erzeugendensystem von $I$. Schreibe die obigen Erzeuger als Linearkombinationen des konstruierten Erzeugendensystems.




}
{} {}


\inputaufgabegibtloesung
{}
{




Es sei $R$ ein
\definitionsverweis {kommutativer Ring}{}{}
und $A$ eine kommutative
$R$-\definitionsverweis {Algebra}{}{.}
$A$ werde durch die Familie

\mavergleichskette
{\vergleichskette
{ a_i
}
{ \in }{A
}
{  }{
}
{    }{
}
{   }{
}
}
{}{}{,}

\mavergleichskette
{\vergleichskette
{ i
}
{ \in }{ I
}
{  }{
}
{    }{
}
{   }{
}
}
{}{}{,}
über $R$ erzeugt. Zeige: Wenn $A$
\definitionsverweis {endlich erzeugt}{}{}
ist, dann wird  $A$ auch von einer endlichen Teilfamilie der $a_i$ erzeugt.




}
{} {}


\inputaufgabe
{}
{




Wir betrachten auf- und absteigende Ketten von
\definitionsverweis {affin-alge\-braischen Mengen}{}{}
in
\mathl{{ {\mathbb A}_{  K }^{  n   } }}{} und von
\definitionsverweis {Idealen}{}{}
in
\mathl{K[X_1 , \ldots , X_n]}{.} Zeige die folgenden Aussagen.
\aufzaehlungvierabc{Für einen endlichen Körper wird jede aufsteigende Kette

\mavergleichskettedisp
{\vergleichskette
{ V_0
}
{ \subseteq} { V_1
}
{ \subseteq} { V_2
}
{ \subseteq} { \ldots
}
{ } {
}
}
{}{}{}
von affin-algebraischen Mengen stationär.
}{Für einen unendlichen Körper und

\mavergleichskette
{\vergleichskette
{ n
}
{ \geq }{ 1
}
{  }{
}
{    }{
}
{   }{
}
}
{}{}{}
wird nicht jede aufsteigende Kette

\mavergleichskettedisp
{\vergleichskette
{ V_0
}
{ \subseteq} { V_1
}
{ \subseteq} { V_2
}
{ \subseteq} { \ldots
}
{ } {
}
}
{}{}{}
von affin-algebraischen Mengen stationär.
}{Für
\zusatzklammer {einen beliebigen Körper und} {} {}

\mavergleichskette
{\vergleichskette
{ n
}
{ \geq }{ 1
}
{  }{
}
{    }{
}
{   }{
}
}
{}{}{}
wird nicht jede absteigende Idealkette

\mavergleichskettedisp
{\vergleichskette
{ {\mathfrak a}_0
}
{ \supseteq} { {\mathfrak a}_1
}
{ \supseteq} { {\mathfrak a}_2
}
{ \supseteq} { \ldots
}
{ } {
}
}
{}{}{}
stationär.
}{Für einen unendlichen Körper und

\mavergleichskette
{\vergleichskette
{ n
}
{ \geq }{ 1
}
{  }{
}
{    }{
}
{   }{
}
}
{}{}{}
gibt es echt absteigende Ketten von affin-algebraischen Mengen beliebiger Länge.
}




}
{} {}


\inputaufgabe
{}
{




Zeige, dass die Menge $\R$ der
\definitionsverweis {reellen Zahlen}{}{}
mit der
\definitionsverweis {metrischen Topologie}{}{}
kein
\definitionsverweis {noetherscher}{}{}
\definitionsverweis {topologischer Raum}{}{}
ist.




}
{} {}






\inputaufgabe
{}
{




Es sei $G$ eine
\definitionsverweis {kommutative Gruppe}{}{.}
Zeige, dass $G$  auf genau eine Weise die Struktur eines
$\Z$-\definitionsverweis {Moduls}{}{}
trägt. Kommutative Gruppen und $\Z$-Moduln sind also äquivalente Objekte.




}
{} {}


\inputaufgabe
{}
{




Es seien $R$ und $A$
\definitionsverweis {kommutative Ringe}{}{.}
Zeige, dass $A$ genau dann eine
$R$-\definitionsverweis {Algebra}{}{}
ist, wenn $A$ ein
$R$-\definitionsverweis {Modul}{}{}
ist, für den zusätzlich

\mathdisp {r (ab) =(ra)b \text{ für alle } r \in R,\, a,b \in A} {  }

gilt.




}
{} {}


\inputaufgabegibtloesung
{}
{




Es sei $V$ ein
\definitionsverweis {Modul}{}{}
über dem
\definitionsverweis {kommutativen Ring}{}{}
$R$. Es seien

\mavergleichskette
{\vergleichskette
{ s_1 , \ldots , s_k
}
{ \in }{ R
}
{  }{
}
{    }{
}
{   }{
}
}
{}{}{}
und

\mavergleichskette
{\vergleichskette
{ v_1 , \ldots , v_n
}
{ \in }{ V
}
{  }{
}
{    }{
}
{   }{
}
}
{}{}{.}
Zeige

\mavergleichskettedisp
{\vergleichskette
{ \left(  \sum_{i = 1}^k s_i  \right) \cdot \left(  \sum_{j = 1}^n v_j  \right)
}
{ =} { \sum_{ 1 \leq i \leq k,\, 1 \leq j \leq n }   s_i \cdot v_j
}
{ } {
}
{ } {
}
{ } {
}
}
{}{}{.}




}
{} {}


\inputaufgabe
{}
{




Es sei $R$  ein
\definitionsverweis {kommutativer Ring}{}{,}
\mathkor {} {M} {und} {N} {} zwei
$R$-\definitionsverweis {Moduln}{}{}
und sei
\maabbdisp {\varphi} {M} {N
} {}
ein
\definitionsverweis {Modulhomomorphismus}{}{.} Zeige die folgenden Aussagen.
\aufzaehlungvier{Für einen
$R$-\definitionsverweis {Untermodul}{}{}

\mavergleichskette
{\vergleichskette
{S
}
{ \subseteq }{ M
}
{  }{
}
{    }{
}
{   }{
}
}
{}{}{}
ist auch das
\definitionsverweis {Bild}{}{}

\mathl{\varphi(S)}{} ein Untermodul von $N$.
}{Insbesondere ist das Bild

\mavergleichskette
{\vergleichskette
{ \operatorname{bild}  \varphi
}
{ = }{ \varphi(M)
}
{  }{
}
{    }{
}
{   }{
}
}
{}{}{}
der Abbildung ein Untermodul von $N$.
}{Für einen Untermodul

\mavergleichskette
{\vergleichskette
{ T
}
{ \subseteq }{ N
}
{  }{
}
{    }{
}
{   }{
}
}
{}{}{}
ist das
\definitionsverweis {Urbild}{}{}

\mathl{\varphi^{-1}(T)}{} ein Untermodul von $M$.
}{Insbesondere ist der
\definitionsverweis {Kern}{}{}

\mathl{\varphi^{-1}(0)}{} ein Untermodul von $M$.
}




}
{} {}






  \zwischenueberschrift{Aufgaben zum Abgeben}


\inputaufgabe
{3}
{




Es sei $R$ ein
\definitionsverweis {kommutativer Ring}{}{}
und sei ${\mathfrak a}$ ein
\definitionsverweis {Ideal}{}{}
mit dem
\definitionsverweis {Restklassenring}{}{}

\mavergleichskettedisp
{\vergleichskette
{ S
}
{ =} { R/{\mathfrak a}
}
{ } {
}
{ } {
}
{ } {
}
}
{}{}{.}
Zeige, dass die Ideale von $S$ eindeutig denjenigen Idealen von $R$ entsprechen, die ${\mathfrak a}$ umfassen.




}
{Zeige, dass das Gleiche für Primideale, Radikalideale und maximale Ideale gilt.} {}


\inputaufgabe
{4}
{




Es sei $R$ ein
\definitionsverweis {noetherscher}{}{}
\definitionsverweis {Integritätsbereich}{}{.}
Zeige, dass sich jedes Element aus $R$ als ein Produkt von
\definitionsverweis {irreduziblen Elementen}{}{}
schreiben lässt.




}
{} {}


\inputaufgabe
{4}
{




Zeige, dass $\mathbb Q$ keine
\definitionsverweis {Algebra von endlichem Typ}{}{} über $\mathbb Z$ ist.




}
{} {}


\inputaufgabe
{4}
{




Es sei $K$ ein
\definitionsverweis {Körper}{}{}
und sei

\mavergleichskette
{\vergleichskette
{ A
}
{ = }{ K[X,Y]
}
{  }{
}
{    }{
}
{   }{
}
}
{}{}{.}
Finde eine
$K$-\definitionsverweis {Unteralgebra}{}{}
von $A$, die nicht endlich erzeugt ist.




}
{} {}


\inputaufgabe
{4}
{




Bestimme zum Ideal

\mavergleichskettedisp
{\vergleichskette
{ I
}
{ =} { \left(  10,6x^2+8,4x^3-12  \right)
}
{ } {
}
{ } {
}
{ } {
}
}
{}{}{}
in $\Z[x]$ die im Beweis zum

Hilbertschen Basissatz

konstruierte Idealkette und das zugehörige Erzeugendensystem von $I$. Schreibe die obigen Erzeuger als Linearkombinationen mit dem konstruierten Erzeugendensystem.




}
{} {}
